\subsection{Counting conventions}
\label{sec:conventions}

To avoid mixing units we count four quantities separately. A \emph{PE slot}
is one PE during one array step; a call with $S$ steps offers $RCS$ slots.
A \emph{useful PE operation} is one PE instruction whose result contributes
to the answer: one addition per output for vector addition and scan, one
MAC per element for the dot product, one multiplication per matrix entry
for matrix--vector and convolution. \emph{Scalar operations} count the
multiplications and additions of the scalar reference (a MAC is two), and
\emph{host operations} are additions performed by the host, such as the
column sums of matrix--vector. Throughout, the useful fraction
\begin{equation}
  \eta=\frac{\text{useful PE operations}}{RC\,S}
  \label{eq:eta}
\end{equation}
is distinct from the \emph{static occupancy}, the fraction of PEs
configured with an arithmetic opcode. For one $4\times4$ diagonal chunk,
four of sixteen PEs are arithmetic (occupancy $1/4$) but four useful sums
take four steps, so $\eta=4/(16\cdot4)=1/16$.

The abstract fabric ceiling is $RC\,f_{\mathrm{clk}}/D$ PE operations per
second, $800\times10^{6}$ for $4\times4$, $D=2$ and 100\,MHz. A mapping
attains at most $\eta$ times this ceiling, i.e.\ $50\times10^{6}$ useful
operations per second for the diagonal, dot and scan mappings and
$200\times10^{6}$ multiplications per second for matrix--vector, before any
communication. The resident dot product, at $D+3=5$ clocks per element,
attains at most $20\times10^{6}$ MACs per second.
